\section{Introduction}
\label{sec:introduction}

\subsection{Problem context}

Generative image models have made it straightforward to produce visually rich
concept art from a short text prompt. That capability does not, by itself, solve
the problem of designing a texture for an existing three-dimensional game asset.
The visible surface of a weapon is stored as a set of UV charts: pieces of the
surface are flattened, rotated, scaled, and packed into a square image. A source
illustration that looks coherent in two dimensions can therefore be split across
unrelated atlas regions, disappear on narrow components, or become inconsistent
between the left, right, and top views. Export requirements, paintable regions,
texture resolution, and seam behavior add further constraints.

This distinction matters for both research and practice. If the generated source
image is treated as the result, an evaluation can reward an attractive rectangle
while overlooking the failures introduced by mapping. Conversely, a simple
tileable pattern may survive arbitrary UV fragmentation but provide little control
over whole-weapon composition. A useful system must connect creative generation to
the actual asset contract and must make its decisions observable.

\skinsmith\ addresses this integration problem. It is a constraint-aware
generative Agent for weapon-skin design, evaluated through a Counter-Strike 2
case study. The system accepts a registered weapon and a compact theme such as
\emph{dragon}. It expands the theme into a visual world, offers several textual
directions, generates several landscape artworks for the selected direction, and
maps every artwork to the target geometry before asking the user to choose. The
chosen source is then processed at full resolution, evaluated, optionally refined
once, and exported with prompts, model identifiers, hashes, metrics, and decisions.

\subsection{Scope}

The research object is the Agent workflow and its geometry-aware texture pipeline,
not a universal Counter-Strike skin generator. The AK-47 is the complete case
study. An M4A4 adapter is used to examine transfer, but it is not presented as an
equivalent production result. The system preserves each asset's existing UV
layout; it does not perform automatic UV unwrapping. Its software renderer provides
controlled technical views rather than photorealistic game-engine rendering.
Steam Workshop publication is outside the project scope.

The group originally registered under the umbrella topic
\emph{Latent Vision $\rightarrow$ Image Generation}. Its initial brief proposed a
broad investigation of image-to-image generation, diffusion models, prompt
engineering, conditioning, and creative outputs such as illustrations, game
concept art, and posters. With fewer than three weeks remaining, the supervisor
advised the group to narrow the work to the technical development of a completed
image-processing Agent and explicitly stated that it did not need to remain the
original broad Latent Vision project. \skinsmith\ is the resulting scoped research
project: it retains image generation and human-AI collaboration from the registered
topic, but concentrates them on one technically testable asset pipeline.

This approved narrowing made it possible to evaluate the complete path from intent
to deployable texture, including negative results, UV failures, selection errors,
and rollback behavior that would be invisible in a source-image-only study. The
registered topic is therefore retained on the title page for administrative
continuity, while the report title names the actual system and research question.

\subsection{Research questions}

The study addresses five questions:

\begin{description}[leftmargin=1.2cm,style=nextline]
  \item[RQ1] Does composing artwork in a mesh-derived continuous weapon space and
  then baking it to the authoritative UV atlas improve mapped asset-seam and
  multi-view performance relative to a generic tileable-pattern baseline?

  \item[RQ2] What failures become observable when candidates are evaluated after
  mapping to the real weapon from left, right, and top views rather than only as
  source images or flat texture atlases?

  \item[RQ3] Can hard-constraint filtering, Pareto-based selection, and a bounded
  correction with rollback prevent an Agent from accepting technically invalid or
  lower-quality updates that a soft aggregate score might prefer?

  \item[RQ4] Can three explicit human checkpoints separate theme intent,
  art-direction choice, and mapped-artwork selection while allowing the Agent to
  begin from one short keyword?

  \item[RQ5] Which parts of the pipeline transfer to a second weapon through a
  \texttt{GameAssetAdapter}, and which geometry, material, camera, and deployment
  contracts remain asset-specific?
\end{description}

\subsection{Contributions}

The project makes five bounded contributions:

\begin{enumerate}
  \item A weapon-space-to-UV generation route that composes selected artwork in a
  canonical frame derived from the target mesh and bakes it through per-UV-pixel
  geometry maps into the existing atlas.
  \item Render-in-the-loop evaluation in which source images, true asset seams,
  and fixed left, right, and top renders are retained as separate evidence.
  \item A small-pool multi-objective policy that separates feasibility from
  preference and permits one local correction with an explicit rollback gate.
  \item A three-checkpoint human-Agent interaction that preserves human authority
  over the theme, direction, and mapped artwork.
  \item An adapter boundary that distinguishes reusable planning and evaluation
  logic from asset-specific mesh, UV, camera, and export contracts.
\end{enumerate}

These claims concern the integration and evaluation of a deployable workflow.
They do not claim the first use of diffusion for mesh texturing, the first
multi-view texture method, or general compatibility with every game asset.

\subsection{Report structure}

\Cref{sec:related} positions the work against diffusion, texture synthesis,
3D-aware texturing, constrained selection, and human-AI systems.
\Cref{sec:method} describes the Agent, asset representation, weapon-space bake,
evaluation, and refinement rules. \Cref{sec:experiments} defines the controlled
experiments and evidence protocol. \Cref{sec:results} reports findings in the
order of the research questions. \Cref{sec:discussion} considers validity,
limitations, project management, and responsible use before
\cref{sec:conclusion} concludes.
